% =====================================================================
%  SAFE INFERENCE — keynote (20 July 2026)
%  SKELETON / OUTLINE, not a finished deck. Frame titles + terse content
%  + figure/reuse hooks, mapped 1:1 to the timed arc in
%  _keynote/safe-inference-notes.md §5. No polish yet — react, then draft.
%
%  Preamble is a compact self-contained copy of ../../slides/day1.tex so
%  this compiles standalone. Reusable frames referenced in comments live
%  in ../../slides/custom/ (motivating_cases.tex, theses.tex).
% =====================================================================
\documentclass[xcolor=dvipsnames]{beamer}

\usepackage{xcolor}
\usepackage{graphicx}
\usepackage{amsmath,amssymb,amsfonts,bm}
\usepackage[english]{babel}
\usepackage{times}
\usepackage[T1]{fontenc}
\usepackage{helvet}
\usepackage{ulem}

\usetheme{Warsaw}
\definecolor{someblue}{rgb}{0,.1,0.8}
\newcommand\makebiggaps{\addtolength{\itemsep}{.5\baselineskip}}
\newcommand\YY{Y(0)}                    % counterfactual shorthand for the talk
\newcommand\E{\mathbb{E}}
\newcommand\ATT{\text{ATT}}
\newcommand\att{\widehat{\text{ATT}}}
\setbeamertemplate{navigation symbols}{}

\title[Safe Inference]{Safe Inference}
\subtitle{Say only what you can defend --- and refuse to abandon what you can}
\author{Chad Hazlett}
\date{20 July 2026}

\begin{document}

\begin{frame}[plain]\titlepage\end{frame}

% =====================================================================
\section{Where we left off}
% ---- BRIDGE from the 2-day course; audience has the machinery fresh ----

\begin{frame}{You left with four theses}
\small
% REUSE: ../../slides/custom/theses.tex  ->  \allTheses  (or thesisThree/Four)
\begin{itemize}\makebiggaps
\item<1-> \textbf{(3)} There are no ``causal methods.'' Credibility comes from the
      \emph{assumptions and the story behind them}, not the technique.
\item<2-> \textbf{(4)} Experiments are not a separate \emph{kind} of knowledge ---
      every design leans on assumptions; the trial just makes one hold by design.
\end{itemize}
\medskip
\onslide<3->{\textbf{This talk is what theses 3 and 4 demand you actually \emph{do}.}}
\end{frame}

% =====================================================================
\section{The near-miss}   % ARC 1 — cold open (~3 min). GBM, unresolved.

\begin{frame}{A treatment we are nearly sure worked}
% TODO: tell it as a person, not a method. GBM. ~30 sec, human.
\begin{itemize}\makebiggaps
\item A therapy. A patient population whose untreated course we know all too well.
\item Strong reason to believe it helped.
\item \textbf{Not credited} --- because it didn't come from an experiment.
\end{itemize}
\vfill
\onslide<2->{\centering\emph{We'll come back to this. (Hold the name.)}}
\end{frame}

\begin{frame}{The question under the near-miss}
\centering\large
How many \emph{defensible} truths do we throw away ---\\[.4em]
and would we even \emph{notice}?
\end{frame}

% =====================================================================
\section{Two errors}      % ARC 2 (~4 min)

\begin{frame}{Two ways to be wrong}
\begin{itemize}\makebiggaps
\item<1-> \textbf{Commission} --- assert an effect that isn't defensible.
\item<2-> \textbf{Omission} --- fail to assert an effect that \emph{is} defensible.
\end{itemize}
\vfill
% NB: evoke Type I / II, do NOT label them — room will hear the ordinary ones.
\onslide<3->{\textbf{We police the first. We are structurally blind to the second.}}
\end{frame}

\begin{frame}{Both live in \emph{your} work, not just medicine}
\begin{itemize}\makebiggaps
\item<1-> Commission: the observational estimate you report with a straight face.
\item<2-> Omission: the effect you \emph{could} defend, dropped the moment a
      referee says ``but it's not an experiment.''
\end{itemize}
\end{frame}

% =====================================================================
\section{The villain}     % ARC 3 (~7 min)

\begin{frame}{One failure, two faces}
\centering\large
A \emph{procedure} sitting in the seat where \emph{judgment} belongs.\\[.6em]
\onslide<2->{\normalsize (Thesis 3, turned into a diagnosis.)}
\end{frame}

\begin{frame}{Face 1 --- the regression tell \hfill\small(commission; home turf)}
\small
\begin{itemize}\makebiggaps
\item Your output looks identical whether your assumptions were mild or heroic.
\item A tight CI comes out \emph{either way}. False precision is the house style.
\end{itemize}
\medskip
% REUSE: ../../slides/custom/motivating_cases.tex — contact/prejudice + microfinance
\onslide<2->{\footnotesize Cases: intergroup contact \& prejudice; microfinance ---
observational stories that RCTs later cut down to size.}
\end{frame}

\begin{frame}{Face 2 --- the RCT heuristic \hfill\small(omission)}
\begin{itemize}\makebiggaps
\item ``Did it come from an experiment?'' treated as the \emph{whole} test of credibility.
\item So recoverable knowledge is discarded because it arrived by another route.
\end{itemize}
\vfill
% Keep two-sided: the enemy straddles the RCT/observational line. Not war on RCTs.
\onslide<2->{\centering\emph{Same failure. Opposite direction.}}
\end{frame}

% =====================================================================
\section{The creed}       % ARC 4 (~3 min)

\begin{frame}{Safe inference}
\begin{enumerate}\makebiggaps
\item<1-> \textbf{Never assert the indefensible.}
\item<2-> \textbf{Never refuse to assert the defensible.}
\end{enumerate}
\vfill
\onslide<3->{\footnotesize Calibrated to the defensible --- no bolder, no more timid.
An idea the lab already argues about; not a slogan coined for this talk.}
\end{frame}

% =====================================================================
\section{The engine}      % ARC 5 (~13 min)

\begin{frame}{One generative question}
\centering\large
Where can we defensibly \textbf{know}, \textbf{bound}, or \textbf{model}
the counterfactual $\YY$ --- without an experiment?
\end{frame}

% ---- METHODS SPINE: they just did 2 days of POs/DAGs. Ground it in notation. ----
\begin{frame}{The estimand, in your notation}
\small
\[ \ATT \;=\; \E\!\left[\,Y_i(1) - Y_i(0)\;\middle|\; D_i = 1\,\right] \]
\begin{itemize}\makebiggaps
\item For a treated unit: $Y_i(1)$ is \emph{observed}; $Y_i(0)$ is the \emph{missing}
      counterfactual.
\item Estimator: $\displaystyle \att = \frac{1}{n_1}\sum_{i:\,D_i=1}
      \big[\,Y_i - \widehat{Y_i(0)}\,\big].$
\end{itemize}
\vfill
\onslide<2->{\textbf{Every design is just a claim about $\widehat{Y_i(0)}$.}
(Thesis 4.) Everything else is bookkeeping.}
\end{frame}

\begin{frame}{Safe inference in one line}
\centering\large
Declare a defensible interval for $\delta$ (the residual bias),\\[.3em]
then report the \emph{entire} ATT band across that interval.
\vfill
\onslide<2->{\normalsize Sensitivity analysis, partial identification, and SCQE
are all instances of this move.}
\end{frame}

% ---- The same move, worked mechanically at three levels of knowledge of Y(0) ----
\begin{frame}{Case 1 --- $\YY$ (nearly) known}
\small
\begin{itemize}\makebiggaps
\item Untreated course near-deterministic (e.g.\ uniformly fatal):
      $Y_i(0)\approx y_0$, essentially known.
\item Then $\displaystyle \att = \bar Y_{\text{treated}} - y_0$ --- a
      \emph{single treated arm} suffices.
\item No control group needed: the counterfactual was never in doubt.
\end{itemize}
\vfill
\onslide<2->{This is the GBM / melanoma regime. Randomization would buy nothing.}
\end{frame}

\begin{frame}{Case 2 --- $\YY$ uncertain: carry a $\delta$}
\small
Write the imputation \emph{with its own error}, and propagate it:
\[ \widehat{Y_i(0)} = m_i + \delta
   \qquad\Longrightarrow\qquad
   \ATT(\delta) = \bar Y_{\text{treated}} - \bar m - \delta . \]
\begin{itemize}\makebiggaps
\item $m_i$: best baseline guess. $\delta$: the residual bias you cannot rule out.
\item Defensible $\delta \in [\delta_{\text{lo}},\,\delta_{\text{hi}}]
      \;\Rightarrow\;$ an ATT \emph{band}, not a point.
\item Whole band clears $0$? \textbf{Defensible claim} (condition 2).
      Straddles $0$? Then say so --- that's the honest answer.
\end{itemize}
\end{frame}

\begin{frame}{SCQE --- model the \emph{trend}, not the confounder}
\small
Cohorts as treatment becomes available; $\pi_c=$ treated share in cohort $c$:
\[ \bar Y_c \;=\; \underbrace{\bar Y(0)_c}_{\text{baseline}} + \ATT\cdot\pi_c
   \;\;\Longrightarrow\;\;
   \Delta\bar Y \;=\; \underbrace{\Delta\bar Y(0)}_{\equiv\,\delta\ (\text{untestable})}
   + \ATT\cdot\Delta\pi . \]
\[ \boxed{\;\ATT(\delta) \;=\; \dfrac{\Delta\bar Y - \delta}{\Delta\pi}\;} \]
\begin{itemize}
\item Ignorability (``no confounding'') is \emph{replaced} by a nameable baseline
      trend $\delta$.
\item Name a defensible interval for $\delta$ $\Rightarrow$ report the ATT band ---
      \emph{the same move as Case 2}.
\end{itemize}
% TODO: reconcile notation with Chad's SCQE slides (Hazlett 2020, Stat.\ Med.);
% he can supply the canonical figures/derivation.
\end{frame}

\begin{frame}{Four ways at $\YY$}
\begin{itemize}\makebiggaps
\item \textbf{Know it} --- prognosis near-deterministic; one treated arm nearly as
      informative as an RCT.
\item \textbf{Bound it} --- partial identification: a range you believe.
\item \textbf{Model its trend} --- SCQE: a nameable, boundable baseline trend
      replaces untestable ``no confounding.''
\item \textbf{Stress-test it} --- sensitivity analysis: how wrong would you have to be?
\end{itemize}
\end{frame}

\begin{frame}{You already believe in this}
\small
% REUSE: ../../slides/custom/motivating_cases.tex — smoking & lung cancer bullet
\begin{itemize}\makebiggaps
\item Smoking causes lung cancer --- among the most trusted claims in science.
\item No RCT. Ever. We believe it because someone formalized \emph{what it would
      take} to be wrong --- and the confounding required was absurd.
\end{itemize}
\vfill
\onslide<2->{\textbf{That is sensitivity analysis. That is safe inference.} You just
never called it that.}
\end{frame}

\begin{frame}{Graceful degradation --- a dimmer, not a switch}
\begin{itemize}\makebiggaps
\item RCT fundamentalism: RCT $\Rightarrow$ believe; else $\Rightarrow$ silence.
\item As knowledge of $\YY$ weakens, the claim weakens \emph{smoothly} through a
      zone of defensible ranges.
\end{itemize}
% TODO FIGURE: switch vs. dimmer schematic.
\end{frame}

\begin{frame}{The deep point}
\centering\large
An RCT's value is \emph{proportional to our ignorance of} $\YY$.\\[.6em]
\onslide<2->{\normalsize Where $\YY$ is knowable by other means, randomization is
redundant --- demanding the trial anyway is ritual, not rigor.}
\end{frame}

\begin{frame}{But: \emph{two} conditions --- this is the sharp part}
\begin{enumerate}\makebiggaps
\item<1-> \textbf{Recoverable $\YY$} --- known or defensibly boundable (narrow $\delta$).
\item<2-> \textbf{Signal clears the band} --- the effect is large relative to the
      $\delta$-band width $+$ selection.
\end{enumerate}
\vfill
\onslide<3->{\footnotesize When a knowable-but-noisy $\YY$ meets a \emph{small}
effect, the band straddles zero --- and \emph{that} is exactly when the RCT is
still doing real work. Match the instrument to where $\YY$ sits.}
\end{frame}

% =====================================================================
\section{The ladder}      % ARC 6 (~14 min) — same engine, escalating, both verdicts

\begin{frame}{TB / isoniazid --- \textsc{restraint} \hfill\small(commission)}
\small
\begin{itemize}\makebiggaps
\item Published SCQE application (Statistics in Medicine, 2020).
\item Naive treated-vs-untreated gap is \emph{dramatic and misleading} --- much
      therapy went to \emph{low-risk} patients.
\item SCQE: state the defensible baseline trend, report the ATT band. Rein it in.
\end{itemize}
% TODO FIGURE: raw gap vs. SCQE ATT band.
\end{frame}

\begin{frame}{GBM --- \textsc{rescue} \hfill\small(omission; the flagship)}
\small
% RESOLVE THE COLD OPEN HERE.
\begin{itemize}\makebiggaps
\item The near-miss from the opening --- named at last.
\item Not naive $\YY$ imputation: a compromised design, sample \emph{mostly but not
      entirely} treated.
\item Use the \emph{whole} study group $+$ SCQE to get $\YY$ for the treated and
      recover the \textbf{ATT}. Safe inference under non-random selection.
\end{itemize}
% TODO FIGURE: study-group -> Y(0) for treated -> ATT band.
\end{frame}

\begin{frame}{Three COVID treatments --- \textsc{scale} \hfill\small(omission, under time pressure)}
\begin{itemize}\makebiggaps
\item Three informative claims impossible under the RCT heuristic.
\item Waiting for the trial was \emph{itself} the costly choice.
\end{itemize}
% TODO: one crisp line per treatment from the dissertation chapters (§7 open item).
\end{frame}

\begin{frame}{A class, not three anecdotes}
\small
\begin{itemize}\makebiggaps
\item \textbf{Melanoma} (vemurafenib): \emph{shares} the structure. Grim, near-fixed
      $\YY$; huge signal. The field even argued the RCT was \emph{unethical}.
\item \textbf{Pancreatic} (daraxonrasib): \emph{boundary case}. $\YY$ equally
      recoverable, but the effect is incremental --- condition 2 \emph{fails}, so
      the confirmatory RCT was warranted.
\end{itemize}
\vfill
\onslide<2->{\centering\emph{The boundary is where the method proves it's honest.}}
\end{frame}

% =====================================================================
\section{The symmetry}    % ARC 7 (~4 min) — the disarming slide

\begin{frame}{Same engine, opposite verdicts}
\centering
\begin{tabular}{ll}
TB / melanoma-\emph{restraint} & \textsc{rein in} an over-claim\\[.4em]
GBM / melanoma-\emph{rescue}   & \textsc{recover} a defensible one\\
\end{tabular}
\vfill
\onslide<2->{\textbf{This is calibration --- not skepticism, not boldness.}}
\end{frame}

\begin{frame}{And it's not just medicine}
\small
\begin{itemize}\makebiggaps
\item \textbf{Rescue:} naloxone / harm-reduction laws --- $\YY$ (overdose-mortality
      trend) boundable, effect can clear the band, RCT unthinkable.
\item \textbf{Restraint:} minimum wage, right-to-carry --- $\YY$ boundable, but the
      effect is small vs.\ baseline volatility. Honest humility, not a bold claim.
\item \textbf{Already done:} pet ownership \& wellbeing (SCQE on panel data) ---
      the discipline in action.
\end{itemize}
\end{frame}

% =====================================================================
\section{The cost}        % ARC 8 (~2 min) — ONE sharp beat, then step off.

\begin{frame}[plain]
\centering\Large
The omission error has a body count too.\\[.6em]
\normalsize In social science it just doesn't leave a body to point to.
\end{frame}

% =====================================================================
\section{The call}        % ARC 9 (~5 min)

\begin{frame}{Norms to pin above a desk}
\small
\begin{enumerate}\makebiggaps
\item State the untestable assumption \emph{before} the estimate, not in the appendix.
\item Report how wrong you'd have to be, to be wrong.
\item Never let a tight interval stand in for a strong assumption.
\item A range you believe beats a point you don't.
\item Do not discard defensible evidence for want of a trial.
\item Reward the honest ``we can't tell yet'' over the confident wrong answer.
\end{enumerate}
\end{frame}

\begin{frame}{What ``rigor'' should mean}
\centering\large
Not the \emph{tightest} interval.\\[.4em]
The \emph{truest} one.
\end{frame}

\begin{frame}[plain]
\centering\large
% Flips the SICSS title "Causality you can believe in?"
Causality you can believe in isn't the causality that sounds most certain ---\\[.5em]
it's the causality still standing after you've told the whole truth\\
about what you assumed.
\end{frame}

\end{document}
